%--- iliad ---
% title: Agent Foundations
% summary: >-
%   Superintelligence alignment may allow only one critical try,
%   with no trial and error on a system more capable than anything yet observed.
%   Agent foundations seeks research directions that survive extreme optimization
%   and self-modification.
%--- end ---
\documentclass[11pt]{article}
\IfFileExists{iliad.sty}{\usepackage[boxes]{iliad}}{\usepackage[boxes]{../iliad}}

% Notation from the source documents, unchanged. The exercise sheet writes its
% own symbols out longhand and needs only \argmin; the rest are the solutions
% document's macros, kept so its proofs are copied byte-for-byte.
%
% Dropped from the sheet's preamble: mathabx (it restyles every math glyph and
% is not in CI's TeX Live; no symbol from it is used), algpseudocode and
% mathtools (nothing in the sheet uses either), and geometry/inputenc/amsmath,
% which iliad.sty supplies.
\DeclareMathOperator*{\argmin}{argmin}
\newcommand{\proves}{\vdash}
\newcommand{\Halts}{\mathsf{Halts}}
\newcommand{\PS}{\mathsf{ProofSeeker}}
\newcommand{\FB}{\mathsf{FairBot}}
\newcommand{\bx}{\Box}
\newcommand{\rd}{r(\delta)}
\newcommand{\DKL}{D_{\mathrm{KL}}}
\DeclareMathOperator{\MI}{MI}
\DeclareMathOperator{\EU}{EU}
\newcommand{\doo}[1]{\mathrm{do}(#1)}
\newcommand{\soln}{\medskip\noindent\textit{Solution.}\ }
% The solutions end their proofs with amsthm's \qed. Spelling out what the
% tombstone is lets the converter carry it to the web (it reads \renewcommand),
% and leaves the PDF marker exactly as it was.
\renewcommand{\qed}{\ensuremath{\square}}
% A Unicode glyph the exercise sheet uses verbatim: the thermostat holding a room
% at "20 ∘ C". inputenc has no mapping for U+2218, so it is declared rather than
% edited out of the author's sentence (same approach as tex/training-dynamics).
\usepackage{newunicodechar}
\newunicodechar{∘}{\ensuremath{^\circ}}

\title{Agent Foundations}
\author{
  \authorname{Daniel C}
  \and
  \authorname{Satya Benson}\\
  \affiliation{Williams College}
}

\begin{document}
\maketitle

\begin{learningoutcomes}
\textbf{(What.)} By the end of the day, students can explain \emph{embedded agency}: why an agent built into the world it acts on (made of the same stuff, smaller than its environment, with no clean input/output boundary, able to model and modify itself) breaks the standard dualistic picture. For each of the day's core research directions (called 'strands') they can state the central problem and a plain-language description of a main result: the complete class theorem (consequentialist foundations), Löb's theorem and the tiling obstacle (self-modification), the logical-induction criterion (logical uncertainty), optimization as local entropy reduction (optimization and thermodynamics), and selection theorems (descriptive agent foundations). They can describe at last one strand in mathematical detail. They can explain how each strand is an offshoot of the central threads \emph{reflective stability} and \emph{embedded agency}.

\textbf{(Why.)} Aligning a superintelligence is a problem we may have to get right on the first critical try, reasoning about a system more capable than anything we have observed, before it exists. That rules out pure trial and error and demands concepts that stay meaningful under extreme optimization pressure and self-modification. Agent foundations supplies (or tries to supply) those concepts, and this day gives students a map of some main research directions and how they fit together.

\textbf{(How.)} Each student goes deep on one pre-reading track beforehand. A one-hour morning lecture lays out the main threads (reflective stability and embedded agency) and previews every strand. Students then read the shared fundamental readings and, in small groups mixing different tracks, run \emph{cross-pollination} discussions that connect their topics to embedded agency. The afternoon makes a few results rigorous through exercises.
\end{learningoutcomes}

\section{Prerequisites}
\begin{itemize}
\item
  Comfort with elementary discrete probability (random variables, expectation, conditional probability).
\item
  Basic formal logic (provability, quantifiers) and basic computability (programs, halting); needed for the Löb and logical-induction strands.
\item
  Basic information theory (entropy, mutual information); needed for the optimization strand.
\item
  No prior agent-foundations background is assumed; every term used on the day is introduced on the day or in the readings. An introductory overview of the AI alignment problem (as in the intro-alignment module) is helpful context but not required.
\item
  Provided refresher notes cover the gaps: a \emph{probability theory} refresher, a \emph{formal logic} refresher, a \emph{computability theory} refresher, and an \emph{information theory, causality, and statistical mechanics} refresher. Point students at the one matching their weakest area before the day.
\end{itemize}

\section{Content}
Exercise sheet: \cref{sec:exercises} below, with worked solutions. The lecture deck is built from this module's own \texttt{slides.tex} and linked at the top of the page.

\subsection{Fast-track}
To get the core in about an hour, or to catch up after missing the day:

\begin{itemize}
\item
  Read the ``Morning lecture'' summary below for the unifying spine (reflective stability, embedded agency, the modelling-vs-implementation split).
\item
  Read \href{https://www.lesswrong.com/posts/i3BTagvt3HbPMx6PN/embedded-agency-full-text-version}{Embedded agency} up to and including section 3.3, then section 4.1, and \href{https://www.lesswrong.com/posts/FWvzwCDRgcjb9sigb/why-agent-foundations-an-overly-abstract-explanation}{Why agent foundations} in full.
\item
  Pick the one strand closest to your background from ``The five strands'' below and read its first (most conceptual) listed reading.
\item
  Skim the takeaways at the end of each strand summary.
\end{itemize}

\subsection{Main content}
Agent foundations is not one theory but several research strands (research directions), each attacking a different facet of the same problem: how to reason about, and build, \emph{stable} agents that are \emph{embedded} in the world they act on. This day is organized around five such strands (consequentialist foundations; Löb's theorem and tiling agents; logical induction; optimization and thermodynamics; descriptive agent foundations). Each student pre-reads one strand in depth before the day; the morning lecture supplies the spine that connects them; the shared fundamental readings and a small-group discussion knit them into a single picture; and an exercise session makes a few of the results rigorous.

Read top to bottom, this section runs: the unifying threads (the morning lecture), then the fundamental readings that give a bird's-eye view, then the five strands each with its readings, then the cross-pollination discussion, the exercises, and a self-check.

\subsubsection{Morning lecture}
\begin{itemize}
\item
  \textbf{Why agent foundations.} We may need key safety properties to hold on the first critical try, for an agent far more capable than any we can study directly. Agent foundations looks for the properties capable agents share \emph{in general}, so we can reason about them in advance.
\item
  \textbf{Reflective stability.} A safety property only matters if it survives the agent's own self-modification: a capable agent may rewrite its code, build a more capable successor, or revise its world model. A property is \emph{reflectively stable} if it is invariant under all of these. Working assumption: under enough optimization pressure, a property persists only if there is some reason it must. Reflective stability is the thread running through every strand.
\item
  \textbf{Robust concepts (``true names'').} Goodhart's law says a proxy that becomes a target stops measuring what we want. So a theory of agency must be built from concepts that do not break under optimization, the ``true names'' of optimization, goals, world models, and embeddedness.
\item
  \textbf{Two pathways of impact (following C. Wyeth).} \emph{Modelling}: build an abstract mathematical model of an idealised capable agent and use it to show why a given alignment proposal would fail. \emph{Implementation}: develop a theory far enough to build and inspect an actual system, often favouring a modular architecture (a separate, inspectable world model, planner, and goal) so each part can be checked.
\item
  \textbf{Dualistic vs embedded agents.} The standard picture (an agent with clean I/O, ``larger than'' and ``outside'' its environment, holding a full world model in its head) is dualistic. A real agent is \emph{embedded}: part of the world, made of the same pieces, smaller than its environment, with no crisp boundary, and able to be copied, modified, or to model itself. Embeddedness is what makes self-improvement, multi-agent reasoning about copies, and self-reference unavoidable, and it is what the rest of the day is about.
\end{itemize}

\subsubsection{Fundamental reading}
\begin{itemize}
\item
  \href{https://www.lesswrong.com/posts/i3BTagvt3HbPMx6PN/embedded-agency-full-text-version}{Embedded agency}: read up to and including section 3.3, then section 4.1. The canonical statement of the embedded-agency problem cluster (the Alexei/Emmy framing the lecture uses).
\item
  \href{https://www.lesswrong.com/posts/FWvzwCDRgcjb9sigb/why-agent-foundations-an-overly-abstract-explanation}{Why agent foundations}: read entirely. Why this abstract, theory-first approach is worth pursuing.
\item
  \href{https://www.lesswrong.com/w/reflectively-consistent-degree-of-freedom}{Reflectively consistent degree of freedom}: read entirely. The precise notion behind a property an agent would not self-modify away, which is the day's recurring theme of reflective stability.
\item
  \href{https://www.lesswrong.com/posts/6mysMAqvo9giHC4iX/what-s-general-purpose-search-and-why-might-we-expect-to-see}{General purpose search}: read entirely. Why a capable mind plausibly contains a retargetable search process.
\end{itemize}

\subsubsection{The five strands (pre-reading tracks)}
\emph{The five strands are different facets of one problem, embedded agency; each student pre-reads one. Every summary below is self-contained at the conceptual level and ends with that strand's readings (start with the first listed, which is the most conceptual). The morning lecture previews all five; the cross-pollination discussion connects them.}

\textbf{1. Consequentialist foundations.} If an agent's preferences are cyclic, an adversary can money-pump it through trades that each look acceptable but leave it strictly worse off. So any agent that reliably avoids such dominated strategies behaves \emph{as if} it maximizes a utility function. The \textbf{complete class theorem} sharpens this: any decision rule that is not dominated (Pareto-optimal across environments) is Bayes-optimal under some prior with full support. This is \emph{representational}, not mechanistic: it says a capable, non-self-defeating agent looks like a Bayesian expected-utility maximizer from the outside, without saying it has a utility function inside, and without telling us \emph{which} utility function (the hard part). \emph{Takeaway: coherence gives a reflectively-stable target (a dominated strategy is one a rational agent would self-modify away from), but it under-determines the agent's actual goals.} \emph{Readings:} \href{https://www.lesswrong.com/s/SgomvxZ3cJWy2SBCu/p/RQpNHSiWaXTvDxt6R}{Coherent decisions imply consistent utilities} (Introduction; ``Why not circular preferences?''; ``Probabilities and expected utilities'' through ``Conditional probability''; Conclusion); \href{https://www.lesswrong.com/posts/73pTioGZKNcfQmvGF/the-measuring-stick-of-utility-problem}{The measuring stick of utility}; \href{https://www.lesswrong.com/posts/sZuw6SGfmZHvcAAEP/complete-class-consequentialist-foundations}{Complete class: consequentialist foundations}.

\textbf{2. Löb's theorem and tiling agents.} A capable agent may build a successor more capable than itself. By \emph{Vingean reflection}, it cannot verify the successor by simulating it (if it could predict the successor's exact moves, it would be that capable already), so it must reason abstractly about the successor's \emph{design}. The natural strategy (``trust the successor because it only takes actions it has proved safe'') needs the parent to trust the successor's proofs. \textbf{Löb's theorem} blocks this: if a consistent system \texttt{L} can prove ``if \texttt{L} proves \texttt{C} then \texttt{C}'', then \texttt{L} already proves \texttt{C}. A consistent system cannot vouch for its own proofs in the abstract, only for a strictly weaker system's. So a naive chain of self-improvements uses an ever-weaker proof system (a ``telomere'' of logical strength that runs out), the \emph{finite descent problem}. The tiling-agents programme studies whether this obstacle can be overcome. \emph{Takeaway: self-trust under self-modification is not free; it runs straight into a logical wall.} \emph{Readings:} \href{https://intelligence.org/files/lob-notes-IAFF.pdf}{Introduction to Löb's theorem} (up to and including section 3); \href{https://www.lesswrong.com/w/vingean-reflection}{Vingean reflection}; \href{https://www.lesswrong.com/posts/QGrX3qK3qxQYK9D4C/walkthrough-of-the-tiling-agents-for-self-modifying-ai-paper}{Walkthrough of the tiling agents paper} (start through ``Finite Descent Problem'', then ``What self-modifying agents need'').

\textbf{3. Logical induction.} Standard Bayesian reasoning assumes \emph{logical omniscience}: the agent instantly knows all consequences of its beliefs. A bounded agent cannot (it may know a program's source yet not its output, or the axioms yet not whether a number is prime). \textbf{Logical induction} (Garrabrant et al.) handles this by picturing a market that prices logical sentences in {[}0,1{]}; the \emph{logical-induction criterion} requires only that no efficient (polynomial-time) trader can exploit the market for unbounded profit, a computable weakening of the Dutch-book argument. That single condition yields convergence and coherence in the limit, timely learning of statistical patterns (it prices ``the nth digit of pi is 7'' near 1/10 without computing it), and \emph{self-trust} (current credence equals a weighted average of expected future credences). \emph{Takeaway: a principled model of how a bounded agent should hold probabilities over facts it has not yet computed.} \emph{Readings:} \href{https://www.lesswrong.com/posts/y5GftLezdozEHdXkL/an-intuitive-guide-to-garrabrant-induction}{An intuitive guide to Garrabrant induction}; \href{https://arxiv.org/pdf/1609.03543}{Logical induction} (chapters 1 and 3; skim chapter 4).

\textbf{4. Optimization and thermodynamics.} A powerful agent reliably steers the world into a narrow region of outcomes, ones extremely unlikely under any random process. This is \emph{local entropy reduction}: concentrating probability mass from a broad initial distribution onto a narrow target. Even a pure predictor has an objective reason to attend to optimizers: naming what an optimizer steers toward predicts the outcome cheaply and robustly, where modelling the initial conditions would be expensive and chaos-fragile. Steering is bounded by information: the \textbf{Touchette-Lloyd} inequality says the entropy reduction a sighted agent achieves over a blind baseline is at most the mutual information between its observations and actions ($\Delta H \leq \Delta H_{\text{blind}}^{\max} + I(X;A)$). \emph{Algorithmic thermodynamics} (Ebtekar and Hutter) replaces ensemble entropy with Kolmogorov complexity, giving laws for individual states and making an embedded agent's knowledge an endogenous physical quantity (algorithmic mutual information between memory and environment), which is exactly its budget for optimization. \emph{Takeaway: optimization is physically constrained, and ``knowing more'' formally means ``being able to optimize more.''} \emph{Readings:} the self-contained note \href{/agency/optimization-thermodynamics}{Optimization and thermodynamics} (read entirely; appendix optional) is the primary reading; its underlying sources are \href{https://www.lesswrong.com/posts/znfkdCoHMANwqc2WE/the-ground-of-optimization-1}{The ground of optimization} (up to and including ``Relationship to Garrabrant and Demski's Embedded Agency''), \href{https://www.lesswrong.com/posts/uKWXktrR7KpbgZAs4/generalized-heat-engine}{Generalized heat engine}, and \href{https://www.lesswrong.com/posts/CJRxQiTKEzior7jGq/algorithmic-thermodynamics-and-three-types-of-optimization}{Algorithmic thermodynamics and three types of optimization}.

\textbf{5. Descriptive agent foundations.} \emph{Normative} agent foundations asks what an ideal agent should look like; \emph{descriptive} asks what agents actually arising in the world (bacteria, neural networks, future AI) look like, and aims to read off their goals, world model, and decision structure from the outside. It works bottom-up from properties of the world (modularity, selection pressures, computational limits). \textbf{Selection theorems} aim to prove results of the form ``any system selected to achieve goal G in environment E must contain structure approximately isomorphic to X'', giving mechanistic rather than merely representational accounts of agency. A key example: the world's \emph{modularity} (it decomposes into sparsely-interacting subsystems) is what makes both world-modelling (Bayesian networks propagate updates locally) and planning (general-purpose search can pursue decoupled subgoals) tractable. \emph{Takeaway: the complementary direction to coherence: not ``non-dominated agents can be described as maximizers'' but ``what pressures make agent-like structure actually arise.''} \emph{Readings:} \href{https://www.lesswrong.com/posts/G2Lne2Fi7Qra5Lbuf/selection-theorems-a-program-for-understanding-agents}{Selection theorems: a program for understanding agents}; \href{https://www.lesswrong.com/posts/TiBsZ9beNqDHEvXt4/how-we-picture-bayesian-agents}{How we picture Bayesian agents}; \href{https://www.lesswrong.com/posts/RuDD3aQWLDSb4eTXP/what-selection-theorems-do-we-expect-want}{What selection theorems do we expect/want}.

\subsubsection{Cross-pollination discussion}
On the day, students who pre-read different strands meet in small mixed groups and work to connect their topics into one picture of embedded agency: how a logical inductor's trust in its future self relates to a tiling agent's trust in its successor, how the coherence account of goals relates to the thermodynamic one, how the representational view of agency (coherence) relates to the mechanistic one (selection theorems). The specific discussion prompts used are listed in the teaching guide.

\subsubsection{Exercise session}
Five exercises (full statements and worked solutions in \cref{sec:exercises}), grouped by topic:

\emph{Logical uncertainty and self-reference.}

\begin{itemize}
\item
  \textbf{Exercise 1: Gödel's second incompleteness theorem} (difficulty 4/5, importance 4/5). Via a self-referential program: a Gödel sentence is true, a consistent system cannot prove its own consistency, and this is exactly the Löbian obstacle to self-trust. Parts (a) true Gödel sentence, (b) no self-consistency proof, (c) the obstacle.
\item
  \textbf{Exercise 2: Löb's theorem} (difficulty 3/5, importance 5/5). The three provability properties (necessitation, distribution, the Löb condition), a full proof of the theorem, and an application: FairBot programs cooperate by Löb's theorem. Parts (a-b) properties, (c-d) the proof, (e) FairBot.
\end{itemize}

\emph{Coherence and consequentialism.}

\begin{itemize}
\item
  \textbf{Exercise 3: The complete class theorem} (difficulty 3/5, importance 5/5). The equivalence between non-dominated strategies and Bayesian expected-utility maximization, via a geometric argument over the convex set of attainable reward vectors. Parts (a) admissibility, (b) faces and difference vectors, (c) constructing the rationalizing prior.
\end{itemize}

\emph{Descriptive agent foundations.}

\begin{itemize}
\item
  \textbf{Exercise 4: The do-divergence theorem} (difficulty 2/5, importance 4/5). Formalizes optimization as outcome concentration and proves that how far an agent can steer outcomes (a KL divergence from the unsteered baseline) is bounded by the mutual information between its observations and actions. The single-step backbone of the Touchette-Lloyd picture.
\item
  \textbf{Exercise 5: Channel additivity} (difficulty 3/5, importance 3/5). Optimal input distributions over independent channels: mutual information decomposes across channels, independence improves throughput, and an optimal policy need not coordinate across a modular environment (connecting modularity to tractable optimization).
\end{itemize}

\subsection{Learn more}
The readings for each strand are linked under that strand in the Main content above. Beyond them:

\textbf{Algorithmic thermodynamics (Ebtekar).} \href{https://arxiv.org/abs/2308.06927}{Foundations of algorithmic thermodynamics}; \href{https://www.mdpi.com/1099-4300/26/9/776}{Modelling the arrows of time with causal multibaker maps}; \href{https://arxiv.org/abs/2408.07818}{Long-time derivation of the Boltzmann equation from hard-sphere dynamics}.

\textbf{Embedded and universal AI (Wyeth).} \href{https://arxiv.org/pdf/2508.16245}{Limit-computable grains of truth}; \href{https://arxiv.org/pdf/2505.17882}{Embeddedness failures in universal artificial intelligence}; \href{https://arxiv.org/pdf/2512.17086}{Value under ignorance}.

\textbf{Other directions.} \href{https://www.lesswrong.com/posts/zB4f7QqKhBHa5b37a/introduction-to-the-infra-bayesianism-sequence}{Introduction to the infra-Bayesianism sequence} (non-realizable environments); \href{https://www.lesswrong.com/posts/ZwshvqiqCvXPsZEct/the-learning-theoretic-agenda-status-2023}{The learning-theoretic agenda}; \href{https://www.lesswrong.com/posts/d2n74bwham8motxyX/optimization-at-a-distance}{Optimization at a distance}; the \href{https://www.lesswrong.com/w/hard-problem-of-corrigibility}{hard problem of corrigibility} and a \href{https://www.lesswrong.com/posts/oLbpfPkdtcknABvvw/the-corrigibility-basin-of-attraction-is-a-misleading-gloss}{critique of the corrigibility basin of attraction}; the original \href{https://intelligence.org/files/TilingAgentsDraft.pdf}{tiling agents draft}; and background notes on \href{https://www2.stat.duke.edu/~pdh10/Teaching/581/LectureNotes/admiss.pdf}{admissibility and the complete class theorem} and the \href{https://www.stat.berkeley.edu/~census/dutchdef.pdf}{Dutch book argument}.

\textbf{Research frontier.} The \emph{agent structure problem} (whether strong optimization provably entails an internal world model), embedded variants of AIXI, and resource-theoretic accounts of instrumental convergence.

\section{Exercises}
\label{sec:exercises}

\begin{quote}
    \textit{The following are exercises on agent foundations. Each problem is broken into a sequence of lemmas leading to a main theorem. \textbf{For each subquestion, try to prove the stated lemma before reading on.} If you get stuck, you may treat the lemma as given and proceed to the next part.}

    \medskip
    \textit{Before diving into a formal derivation, try to build an intuition for \textbf{why} the statement should be true. Even if you don't complete the proof, having a clear intuitive picture of what's going on is more valuable than a mechanical derivation you don't understand. Don't worry if some of the terminology is unfamiliar --- the exercises are designed to be self-contained, and it should be possible to follow the questions from context.}
\end{quote}
\noindent
\textbf{Formal systems and programs.} A \emph{formal system} is a precise set of rules for
deriving mathematical statements from axioms. Fix a formal system $L$ that is powerful
enough to reason about programs (for instance, it can express statements about
arithmetic, and any program can be encoded as a mathematical object that $L$ can talk
about). We write $L \vdash \varphi$ to mean that the statement $\varphi$ is
\emph{provable} in $L$, i.e.\ there exists a finite sequence of steps, each justified by
the rules of $L$, that derives $\varphi$.

\medskip
\noindent
We say $L$ is \emph{consistent} if it never proves a contradiction. We write $\bot$ for
a fixed contradictory statement (such as $0 = 1$), so consistency means $L \nvdash \bot$.
We assume throughout that $L$ is consistent.

\medskip
\noindent
\textbf{Programs.} By a \emph{program} we mean a mechanical procedure that follows a
fixed list of instructions. A program may \emph{halt} (finish and produce an output)
or \emph{run forever} (keep executing without ever stopping). Since $L$ can reason about
programs, it can express the statement ``program $M$ halts'', which we write as
$\mathsf{Halts}(M)$.

\medskip
\noindent
\textbf{The bridge between $L$ and programs.} Formal systems and programs are
intimately connected, and the key to these exercises is switching back and forth between
the two perspectives:

\begin{itemize}
  \item \textbf{From programs to $L$} (concrete outputs become proofs)\textbf{.} If a
    program concretely produces an output (e.g.\ it halts after some number of steps,
    or it finds a string with a certain property), then $L$ can verify this by tracing
    through the execution step by step. In particular: if a program actually halts,
    then $L$ can prove that it halts.

  \item \textbf{From $L$ to programs} (proofs can be found by search)\textbf{.} Proofs
    in $L$ are finite strings that can be checked mechanically. So for any statement
    $P$, we can write a program that searches through all possible strings, checks
    whether each one is a valid $L$-proof of $P$, and halts if it finds one:
    \[
      \texttt{ProofSeeker}(P) \;:=\;
      \text{``try every string; halt iff one is a valid $L$-proof of $P$.''}
    \]
    This program halts if and only if $P$ is provable in $L$.
\end{itemize}

\medskip
\noindent
Whenever you derive a fact about $L$ (e.g.\ that some statement is or isn't provable),
ask what that implies for the corresponding proof-search program, and vice versa.

\medskip
\noindent
\textbf{The provability predicate $\Box P$.} We write $\Box P$ for the statement,
\emph{expressed within $L$ itself}, that ``$P$ is provable in $L$.'' This is a genuine
mathematical statement that $L$ can reason about, because it is equivalent to the claim
that $\texttt{ProofSeeker}(P)$ halts, and $L$ can talk about programs.

\medskip
\noindent
The crucial distinction is between $L \vdash P$ and $L \vdash \Box P$:

\begin{itemize}
  \item $L \vdash P$ means that $P$ is provable: there \emph{exists} a concrete proof
    of $P$ in $L$. This is a fact about $L$ that we observe from the outside.

  \item $L \vdash \Box P$ means that $L$ has proved a statement \emph{about itself}:
    namely, that a proof of $P$ exists (equivalently, that
    $\texttt{ProofSeeker}(P)$ halts). But this is a claim $L$ is making about
    the  $\texttt{ProofSeeker}(P)$ program, not a direct certificate for $P$.
\end{itemize}



\medskip
\noindent
\textbf{Intuition.} To see why $L \vdash P$ and $L \vdash \Box P$ are conceptually
distinct, consider the two different ways $L$ might prove that
$\texttt{ProofSeeker}(P)$ halts. The first is to actually trace through its execution:
if it halts after, say, a million steps, $L$ can verify this step by step, and the proof
that $\texttt{ProofSeeker}(P)$ found along the way is itself a direct $L$-proof of $P$.
In this case, $L \vdash \Box P$ and $L \vdash P$ seems to come hand in hand. But there is a
second way: $L$ might reason \emph{abstractly} about the program's behaviour without
ever simulating it. (This is analogous to how you might argue that a sorting algorithm
must eventually finish without tracing through every swap it makes.) Such a proof
establishes that $\texttt{ProofSeeker}(P)$ halts --- and therefore that \emph{some}
proof of $P$ exists --- but the proof itself is about the \emph{program}, not about
$P$. It need not contain, or even hint at, what the actual proof of $P$ looks like.
This is the gap that $L$ cannot close in general: knowing abstractly that a proof is
``out there'' is not the same as having the proof in hand.\\\\



\medskip
\noindent
\textbf{What ``provable in $L$'' means (and what it does not).} It is
important to distinguish between being \emph{convinced} that something is true and
\emph{proving it in $L$}. When we say ``$L$ can carry out this argument'' or
``$L$ proves $P$,'' we do not mean that a reasonable person reading the argument
would find it convincing. We mean something much more specific: that there exists a
sequence of formulas, each of which is either an axiom of $L$ or follows from
earlier formulas by one of $L$'s explicitly listed inference rules, and whose last
line is $P$. The formal system $L$ is a \emph{machine}: it has no understanding, no
intuition, and no ability to say ``well, this obviously follows.'' Every single step
must be justified by a specific rule.\\

\noindent
From the outside, we can see that if $\texttt{ProofSeeker}(P)$ halts then a proof of $P$
exists, so $P$ is provable. But can $L$ carry out this reasoning internally, always concluding
$P$ from $\Box P$? L\"ob's theorem (Exercise 2) shows that the answer is no: any
consistent $L$ that derives $P$ from $\Box P$ for all $P$ is in fact inconsistent.

\begin{exercise}[G\"odel's second incompleteness theorem]
\label{ex:godel}
\noindent
\textbf{Key fact.} If a program $M$ actually halts (say, after 17 steps), then $L$ can
verify this by checking the execution step by step, so $L \vdash \mathsf{Halts}(M)$.
However, if $M$ runs forever, $L$ cannot necessarily prove
$\neg\mathsf{Halts}(M)$. (Informally: it is easy to certify that something stops, because
you just exhibit the stopping point; but certifying that something runs \emph{forever}
is much harder, because you cannot check infinitely many steps.)

\medskip
\noindent
In fact, no consistent formal system can correctly settle the question ``does $M$ halt?''
for \emph{every} program $M$. To see why: if $L$ could do this, we could write a program
that, given any $M$, searches for an $L$-proof of either $\mathsf{Halts}(M)$ or
$\neg\mathsf{Halts}(M)$. Since $L$ is assumed to settle every case, this search would
always find a proof and halt, giving us a mechanical procedure that decides whether
any program halts. But such a procedure cannot exist (this is the \emph{undecidability
of the halting problem}, a fundamental result in computer science that we take as
given here).

\medskip
\noindent
\textbf{A self-referencing program.} It is possible to write programs that refer to
their own source code. (As a simple example, a program can carry its own source code
as a string and then operate on it.) Using this idea, define the following program:
\[
  Z(A) \;:=\; \text{``search for an $L$-proof of
  $\neg\mathsf{Halts}(A(A))$; halt if one is found.''}
\]
Here $A(A)$ means ``run program $A$ with its own source code as input.'' So $Z(A)$
searches for a proof that the program $A$-run-on-itself runs forever.

\medskip
\noindent
Now consider feeding $Z$ its own source code. The program $Z(Z)$ searches for an
$L$-proof that $Z(Z)$ runs forever. Define the statement:
\[
  G \;:=\; \neg\mathsf{Halts}(Z(Z)).
\]
In words: $G$ says ``$Z(Z)$ runs forever.'' Notice the self-referential
structure: $Z(Z)$ halts if and only if it finds an $L$-proof of $G$ (i.e.\ a proof that
$Z(Z)$ runs forever).

\medskip
\noindent
\textbf{Part 1(a).} Show that $G$ is true, assuming $L$ is consistent.

\medskip
\noindent
\emph{Hint: Consider two cases. Either $Z(Z)$ halts or it doesn't. In each case, use
the bridge between programs and $L$ (if a program halts, $L$ can prove it; if $L$
proves something, the corresponding proof-search program finds that proof and halts)
to derive what follows. One of the two cases leads to a contradiction with the
consistency of $L$.}

\medskip
\noindent
\textbf{Part 1(b).} Show that if $L$ can prove its own consistency (i.e.\
$L \vdash \neg\Box\bot$), then $L$ is in fact inconsistent.

\medskip
\noindent
\emph{Hint: In Part 1(a), you argued from outside $L$ that $G$ is true, and the
argument used only one assumption about $L$: that $L$ is consistent. If $L$ can prove
its own consistency, then every step of your outside argument can be carried out
\textbf{inside} $L$ as a formal derivation. What would $L$ then be able to prove?
And what would the corresponding program do?}

\medskip
\noindent
\textbf{Part 1(c).} Suppose $L$ can vouch for all of its own proofs, meaning
$L \vdash \Box P \to P$ for every statement $P$. (Read this as: ``whenever $L$ can
prove $P$, then $P$ is actually true,'' and $L$ itself asserts this.) Show that $L$ is
inconsistent, in two steps:

\begin{itemize}
  \item[(i)] First, show that $L$ proves it never proves anything false. That is: for
    any $P$ with $L \vdash \neg P$, show that $L \vdash \neg\Box P$.\\
    \emph{Hint: The statement ``if $A$ then $B$'' is logically equivalent to ``if not
    $B$ then not $A$''. Apply this to $\Box P \to P$.}
  \item[(ii)] Apply (i) with $P = \bot$, using the fact that $\neg\bot$ (``a
    contradiction is false'') is a tautology. Conclude that
    $L \vdash \neg\Box\bot$, and use Part 1(b) to finish.
\end{itemize}

\medskip
\noindent
\textbf{Remark (Self-trust, tiling agents, and the L\"obian obstacle).}

\medskip
\noindent
Any sufficiently advanced AI may eventually be able to modify its own code or build a
successor system more capable than itself. But this raises a subtle problem. If the
successor is genuinely smarter, the original agent \emph{cannot} predict exactly what
it will do --- just as the programmers of a chess engine can reason that their program
is ``trying to win'' without knowing its exact moves. So the original agent cannot
verify its successor's safety by simulating it move by move. Instead, it must reason
\emph{abstractly} about the successor's design: ``whatever my successor does, it will
only take actions that it has proved lead to good outcomes.''

\medskip
\noindent
This reasoning strategy is called \emph{tiling}: the parent agent $A_1$ builds a
child agent $A_0$, and wants to conclude not merely that $A_0$ will only take actions
that $A_0$ has \emph{proved to be safe}, but that those actions \emph{actually are}
safe. After all, $A_1$ can verify from $A_0$'s source code that $A_0$ says ``only
take action $x$ if I can prove that $x$ leads to good outcomes.'' But this only tells
$A_1$ that $A_0$ acts on what $A_0$'s proof system certifies --- it does not yet tell
$A_1$ that what $A_0$'s proof system certifies is actually \emph{true}. To close this
gap, $A_1$ needs to be able to prove, within its own reasoning, that $A_0$'s proof
system is \emph{sound}: whenever $A_0$'s system proves $P$, then $P$ is actually
true. When both agents use the same formal system $L$, this amounts to
$L \vdash \Box P \to P$ for all $P$.

\medskip
\noindent
Part 1(c) shows this is impossible: any consistent system that asserts
$\Box P \to P$ for all $P$ is already inconsistent. A consistent system cannot vouch
for its own proofs in the abstract --- it can only trust a proof once it has
\emph{witnessed} it directly. This is the \textbf{L\"obian obstacle}: the barrier
created by L\"ob's theorem to self-trusting formal reasoning.

\medskip
\noindent
One might hope to work around this by having the parent use a \emph{stronger} proof
system than the child: a stronger system \emph{can} trust a weaker one's proofs. But
this means each successive agent in a chain of self-improvements must use a strictly
weaker proof system than its predecessor, resulting in a ``telomere'' of logical
strength that shortens with each generation. Eventually the chain runs out of trust.
The \emph{tiling agents} research programme studies how (and whether) this obstacle can
be overcome, seeking agent architectures that can undergo indefinite self-improvement
without their reasoning guarantees degrading at each step.
\end{exercise}

\begin{solution}[ex:godel]
Recall $Z(A)$ searches for an $L$-proof of $\neg\Halts(A(A))$ and halts if it finds one, and
\[
  G \;:=\; \neg\Halts(Z(Z)).
\]
By construction $Z(Z)$ searches for an $L$-proof of $G$, so
\begin{equation}\label{eq:selfref}
  Z(Z)\text{ halts} \iff L\proves G. \tag{$\star$}
\end{equation}

\paragraph{Part 1(a).} $G$ is true (assuming $L$ consistent).

\soln Two cases.
\begin{itemize}[leftmargin=1.5em]
\item \emph{$Z(Z)$ halts.} By \eqref{eq:selfref}, $L\proves G$, i.e.\ $L\proves\neg\Halts(Z(Z))$. But $Z(Z)$ actually halts, so by the bridge $L\proves\Halts(Z(Z))$. Then $L$ proves both $\Halts(Z(Z))$ and its negation, contradicting consistency.
\item \emph{$Z(Z)$ runs forever.} Then $\neg\Halts(Z(Z))$, i.e.\ $G$, is true.
\end{itemize}
Consistency rules out the first case, so $Z(Z)$ runs forever and $G$ is true. \qed

\paragraph{Part 1(b).} If $L\proves\neg\bx K$ (i.e.\ $L$ proves its own consistency) then $L$ is inconsistent.

\soln The argument of 1(a) used only the consistency of $L$, and each step is a finite manipulation of programs and proofs that $L$ can formalize. Reading it inside $L$: ``$Z(Z)$ halts'' yields $\bx G$ (by construction) and $\bx\Halts(Z(Z))$ (the bridge), hence $\bx K$; contrapositively $L\proves \neg\bx K \to G$. If $L\proves\neg\bx K$, modus ponens gives
\[
  L\proves G,\qquad\text{i.e.}\qquad L\proves\neg\Halts(Z(Z)).
\]
But $L\proves G$ means $Z(Z)$'s search finds a proof of $G$, so $Z(Z)$ halts; by the bridge $L\proves\Halts(Z(Z))$. Now $L$ proves both $\Halts(Z(Z))$ and $\neg\Halts(Z(Z))$, so $L$ is inconsistent. \qed

\paragraph{Part 1(c).} If $L\proves \bx P\to P$ for every $P$, then $L$ is inconsistent.

\soln
\emph{(i)} Fix $P$ with $L\proves\neg P$. The contrapositive of $\bx P\to P$ is $\neg P\to\neg\bx P$, so from $L\proves\bx P\to P$ we get $L\proves\neg P\to\neg\bx P$. With $L\proves\neg P$, modus ponens gives $L\proves\neg\bx P$.

\emph{(ii)} Take $P=K$. Since $\neg K$ is a tautology, $L\proves\neg K$, so by (i) $L\proves\neg\bx K$. Thus $L$ proves its own consistency, and Part 1(b) makes $L$ inconsistent. \qed
\end{solution}

\begin{exercise}[L\"ob's theorem]
\label{ex:lob}
\noindent
L\"ob's theorem says: if $L \vdash \Box C \to C$ (i.e.\ $L$ can prove ``if $C$ is
provable then $C$ is true''), then $L \vdash C$ (i.e.\ $C$ is already provable in
$L$). In other words, the only statements for which $L$ can close the gap between
``provably provable'' and ``provable'' are the ones that were already provable to
begin with.

\medskip
\noindent
The proof uses three properties of the provability predicate $\Box$. We state them
here together with informal explanations of what they say about the proof-search
program $\texttt{ProofSeeker}$.

\medskip
\begin{center}
\begin{tabular}{lll}
  \textbf{(N)} & \emph{Necessitation.} & If $L \vdash \varphi$ then
    $L \vdash \Box\varphi$. \\
  \textbf{(K)} & \emph{Distribution.} & $L \vdash \Box(\varphi \to \psi) \to
    (\Box\varphi \to \Box\psi)$. \\
  \textbf{(4)} & \emph{L\"ob condition.} & $L \vdash \Box\varphi \to
    \Box\Box\varphi$.
\end{tabular}
\end{center}

\medskip
\noindent
\textbf{Part 2(a). Understanding Necessitation.} Necessitation says: if $\varphi$ is
provable in $L$, then $L$ can prove that $\varphi$ is provable. Explain why this is
true from the perspective of $\texttt{ProofSeeker}$.

\medskip
\noindent
\emph{Hint: If a proof of $\varphi$ exists, then $\texttt{ProofSeeker}(\varphi)$ will
find it and halt.}

\medskip
\noindent
\textbf{Part 2(b). Understanding Distribution.} It is helpful to think of a proof of
$\varphi \to \psi$ by analogy with a \emph{function}: given any proof of $\varphi$ as
input, one can mechanically produce a proof of $\psi$ as output (by writing down the
proof of $\varphi$, attaching the proof of $\varphi \to \psi$, and applying the
logical rule that from $\varphi$ and $\varphi \to \psi$ one may conclude $\psi$).

\medskip
\noindent
With this analogy, $\Box(\varphi \to \psi)$ says that $L$ has proved such a
``proof-transforming function'' exists. Distribution then says: if $L$ knows that a
function from $\varphi$-proofs to $\psi$-proofs exists, and $L$ knows that a
$\varphi$-proof exists, then $L$ can conclude that a $\psi$-proof exists.

\medskip
\noindent
Explain why Distribution is true from the perspective of $\texttt{ProofSeeker}$.

\medskip
\noindent
\emph{Hint: If both $\texttt{ProofSeeker}(\varphi \to \psi)$ and
$\texttt{ProofSeeker}(\varphi)$ halt, what can you do with their outputs? How can $L$ deduce that $\texttt{ProofSeeker}(\psi)$ halts?}

\bigskip
\noindent
We now prove L\"ob's theorem. The key ingredient is a self-referential sentence
constructed using the same idea as in Exercise~1 (a sentence that talks about its
own provability). Specifically, there exists a sentence $\lambda$ such that
\[
  L \vdash \lambda \;\leftrightarrow\; (\Box\lambda \to C).
\]
In words: $\lambda$ says ``if I am provable, then $C$ is true.'' (The existence of
such a sentence is guaranteed by the same self-referential construction used to build
$G$ in Exercise~1; we take it as given here.)

\medskip
\noindent
\textbf{Part 2(c).} Show that $L \vdash \Box\lambda \to \Box C$.

\medskip
\noindent
\emph{Hint: The L\"ob sentence says $L \vdash \lambda \to (\Box\lambda \to C)$.
Apply Necessitation to get this fact inside a $\Box$, then use Distribution twice
(once to ``unwrap'' the outer implication, once to handle $\Box\lambda \to C$
inside). Use the L\"ob condition to handle the resulting $\Box\Box\lambda$.}

\medskip
\noindent
\textbf{Part 2(d).} Now assume $L \vdash \Box C \to C$. Using Part 2(c) and the
L\"ob sentence, derive $L \vdash C$.

\medskip
\noindent
\emph{Hint: From Part 2(c), you have $L \vdash \Box\lambda \to \Box C$. Chain this
with the assumption $L \vdash \Box C \to C$ to get
$L \vdash \Box\lambda \to C$. Now compare this with what $\lambda$ says about
itself.}

\medskip
\noindent
\medskip
\noindent
\textbf{Remark (The Santa Claus paradox and what L\"ob adds beyond G\"odel).}

\medskip
\noindent
The L\"obian sentence $\lambda \leftrightarrow (\Box\lambda \to C)$ is a formal
analogue of the \emph{Santa Claus sentence} (also known as Curry's paradox). Consider
the sentence $S$: ``If this sentence is true, then Santa Claus exists.'' Let us try
to determine whether $S$ is true or false.

\medskip
\noindent
Well, $S$ is an ``if \ldots\ then \ldots'' statement, so to check whether it is true,
let us assume the ``if'' part and see whether the ``then'' part follows. So assume $S$
is true. Since $S$ says ``if $S$ is true then Santa Claus exists,'' and we are
assuming $S$ is true, it follows that Santa Claus exists. We have therefore shown: if
$S$ is true, then Santa Claus exists. But that is exactly what $S$ says! So $S$ is
true. And since $S$ is true and $S$ implies Santa Claus exists, Santa Claus exists.
Since nothing about this argument was specific to Santa Claus, the same reasoning
``proves'' any statement whatsoever.

\medskip
\noindent
The reason $L$ does not fall prey to this paradox is that $L$ cannot form a sentence
that refers to its own \emph{truth}. (A fundamental result called Tarski's theorem
shows that no sufficiently powerful formal system can define a truth predicate for
itself.) What $L$ \emph{can} do is refer to its own \emph{provability}: the predicate
$\Box\varphi$ is a legitimate statement within $L$. The L\"obian sentence $\lambda$
therefore substitutes ``provable'' for ``true,'' asserting ``if I am \emph{provable},
then $C$.'' The proof of L\"ob's theorem shows that this substitution is enough to
force $L \vdash C$ --- but only when $\Box C \to C$ is already assumed, not
unconditionally.

\medskip
\noindent
This mirrors a pattern: replacing ``true'' with ``provable'' transforms semantic
paradoxes into precise theorems. The liar's paradox (``this sentence is false'')
becomes G\"odel's sentence (``this sentence is not provable''), yielding the
incompleteness theorems. The Santa Claus paradox (``if this sentence is true, then
$C$'') becomes the L\"obian sentence (``if this sentence is provable, then $C$''),
yielding L\"ob's theorem.

\medskip
\noindent
\textbf{What L\"ob adds beyond G\"odel.} G\"odel's second incompleteness theorem
(Exercise~1) says that $L$ cannot prove its own consistency. This is already a
serious obstacle, but one might hope that consistency is a special case --- perhaps
$L$ can still trust its proofs in less sweeping ways. L\"ob's theorem crushes this
hope completely. It says that for \emph{any} statement $C$, if $L$ can prove ``my
provability of $C$ implies $C$ is true'' (i.e.\ $L \vdash \Box C \to C$), then $C$
was already provable. There are \emph{no} statements, not for which $L$ can assert ``well, if I \emph{could} prove this, it would be
true'' without already being able to prove them. $L$ does not trust its own proofs
until it has witnessed them directly.

\medskip
\noindent
This is what makes L\"ob's theorem, rather than G\"odel's, the fundamental obstacle
for tiling agents. Recall the setup from Exercise~1: a parent agent $A_1$ builds a
child $A_0$ that only takes actions it can prove to be safe. For $A_1$ to trust
$A_0$, it needs to know that $A_0$'s proofs track reality --- that is, $A_1$ needs
$\Box P \to P$ for the statements $P$ that $A_0$ might act on. G\"odel tells us
$A_1$ cannot prove $A_0$'s system is \emph{consistent}. But L\"ob tells us something
far stronger: $A_1$ cannot even trust $A_0$'s system on a \emph{case-by-case} basis.
For any individual statement $C$, the only way $L$ can derive ``if my proof system
proves $C$, then $C$ is really true'' is if $C$ was already provable --- in which
case the trust was never needed in the first place.

\medskip
\noindent
\textbf{Part 2(e).} In a one-shot Prisoner's Dilemma, two players each choose to
either \emph{cooperate} ($C$) or \emph{defect} ($D$). In this variant, instead of
choosing directly, each player submits a \emph{program} that receives the opponent's
source code as input and outputs $C$ or $D$. Consider the following agent,
\emph{FairBot}:

\medskip
\begin{center}
\begin{tabular}{l}
\textbf{algorithm} $\texttt{FairBot}(\text{opponent})$:\\
\quad Search for an $L$-proof that
  $\texttt{opponent}(\texttt{FairBot}) = C$.\\
\quad \textbf{if} proof found \textbf{then return} $C$\\
\quad \textbf{else return} $D$
\end{tabular}
\end{center}

\medskip
\noindent
In words: FairBot cooperates with an opponent if and only if it can find an $L$-proof
that the opponent cooperates with FairBot. Note that FairBot is
\emph{unexploitable}: if $L$ is sound (i.e.\ $L$ only proves true statements), then
FairBot never cooperates with an opponent that defects against it.

\medskip
\noindent
The interesting question is what happens when FairBot plays against itself. At first
glance, both mutual cooperation and mutual defection seem like stable outcomes.

\medskip
\noindent
Consider two copies $\texttt{FairBot}_1$ and $\texttt{FairBot}_2$ (identical programs
with different implementations). Let $A$ be the statement
``$\texttt{FairBot}_1(\texttt{FairBot}_2) = C$'' and $B$ be the statement
``$\texttt{FairBot}_2(\texttt{FairBot}_1) = C$.'' Prove that
$L \vdash A \wedge B$, i.e.\ that the two FairBots mutually cooperate. Use L\"ob's
theorem.

\medskip
\noindent
\emph{Hint: From FairBot's source code, $\Box A$ implies that
$\texttt{FairBot}_2$ finds a proof that $\texttt{FairBot}_1$ cooperates, so
$\texttt{FairBot}_2$ cooperates, i.e.\ $B$ is true. Similarly $\Box B$ implies $A$.
Combine these to show $L \vdash \Box(A \wedge B) \to (A \wedge B)$, and apply L\"ob's
theorem.}
\end{exercise}

\begin{solution}[ex:lob]
We use Necessitation (N) $L\proves\varphi\Rightarrow L\proves\bx\varphi$, Distribution (K) $L\proves\bx(\varphi\to\psi)\to(\bx\varphi\to\bx\psi)$, and the L\"ob condition (4) $L\proves\bx\varphi\to\bx\bx\varphi$.

\paragraph{Part 2(a). Necessitation.} If $L\proves\varphi$ there is a concrete proof of $\varphi$; searching all strings, $\PS(\varphi)$ meets it and halts. ``$\PS(\varphi)$ halts'' is exactly $\bx\varphi$, and a halting run is finite, so by the bridge $L\proves\bx\varphi$. \qed

\paragraph{Part 2(b). Distribution.} A proof of $\varphi\to\psi$ turns a proof of $\varphi$ into one of $\psi$ (concatenate the two proofs and apply modus ponens). So if $\PS(\varphi\to\psi)$ and $\PS(\varphi)$ both halt, their outputs combine into a proof of $\psi$, whence $\PS(\psi)$ halts. This combining is a finite procedure $L$ can carry out, so $L\proves\bx(\varphi\to\psi)\to(\bx\varphi\to\bx\psi)$. \qed

\medskip
The L\"ob sentence $\lambda$ satisfies $L\proves\lambda\leftrightarrow(\bx\lambda\to C)$.

\paragraph{Part 2(c).} $L\proves\bx\lambda\to\bx C$.

\soln From $L\proves\lambda\to(\bx\lambda\to C)$:
\begin{align*}
\text{(N):}\quad & L\proves \bx\!\big(\lambda\to(\bx\lambda\to C)\big),\\
\text{(K):}\quad & L\proves \bx\lambda\to\bx(\bx\lambda\to C), \\
\text{(K):}\quad & L\proves \bx(\bx\lambda\to C)\to(\bx\bx\lambda\to\bx C),
\end{align*}
and chaining the last two, $L\proves\bx\lambda\to(\bx\bx\lambda\to\bx C)$. By (4), $L\proves\bx\lambda\to\bx\bx\lambda$. Propositionally, from $\bx\lambda$ we obtain $\bx\bx\lambda$ and then $\bx C$, so $L\proves\bx\lambda\to\bx C$. \qed

\paragraph{Part 2(d).} Assuming $L\proves\bx C\to C$, derive $L\proves C$.

\soln Chaining $L\proves\bx\lambda\to\bx C$ (2(c)) with $L\proves\bx C\to C$ gives $L\proves\bx\lambda\to C$. The L\"ob equivalence gives the converse direction $L\proves(\bx\lambda\to C)\to\lambda$, so $L\proves\lambda$. By (N), $L\proves\bx\lambda$; with $L\proves\bx\lambda\to C$, modus ponens yields $L\proves C$. \qed

\paragraph{Part 2(e). FairBot.} With $A:=$ ``$\FB_1(\FB_2)=C$'' and $B:=$ ``$\FB_2(\FB_1)=C$'', show $L\proves A\wedge B$.

\soln From the source code, $\FB_1$ returns $C$ iff it finds an $L$-proof that its opponent $\FB_2$ returns $C$ against it, i.e.\ a proof of $B$; reading this off the code,
\[
  L\proves \bx B\to A,\qquad L\proves \bx A\to B .
\]
For any $\varphi$, $\bx(A\wedge B)\to\bx\varphi$ whenever $A\wedge B\to\varphi$ is a tautology: indeed (N) gives $\bx(A\wedge B\to\varphi)$ and (K) then gives $\bx(A\wedge B)\to\bx\varphi$. Applying this to $\varphi=A$ and $\varphi=B$,
\[
  L\proves \bx(A\wedge B)\to\bx A,\qquad L\proves \bx(A\wedge B)\to\bx B.
\]
Hence, assuming $\bx(A\wedge B)$: we get $\bx A$ and $\bx B$, then $B$ (from $\bx A\to B$) and $A$ (from $\bx B\to A$), so $A\wedge B$. Thus
\[
  L\proves \bx(A\wedge B)\to(A\wedge B),
\]
and L\"ob's theorem with $C:=A\wedge B$ gives $L\proves A\wedge B$. \qed
\end{solution}

\begin{exercise}[The Complete Class Theorem]
\label{ex:complete-class}
\noindent
\textbf{The setup: decision-making under uncertainty.} Imagine you must choose an
action, but you don't know which \emph{environment} you are in. There are finitely
many actions $\mathcal{A} = \{a_1, \ldots, a_k\}$ and finitely many possible
environments $\Omega = \{\omega_1, \ldots, \omega_n\}$. If you take action $a$ and the
true environment turns out to be $\omega$, you receive a reward
$u(a, \omega) \in \mathbb{R}$.

\medskip
\noindent
\textbf{Example.} You are deciding whether to carry an umbrella ($a_1$) or not
($a_2$). The environment is either ``rainy'' ($\omega_1$) or ``sunny'' ($\omega_2$).
The rewards might be:
\[
\begin{array}{c|cc}
 & \omega_1 \text{ (rain)} & \omega_2 \text{ (sun)} \\
\hline
a_1 \text{ (umbrella)} & 8 & 5 \\
a_2 \text{ (no umbrella)} & 2 & 10
\end{array}
\]

\medskip
\noindent
\textbf{Decision rules.} A \emph{decision rule}
$\delta = (\lambda_1, \ldots, \lambda_k)$ is a (possibly randomized) strategy: you
play action $a_j$ with probability $\lambda_j$, where $\lambda_j \geq 0$ and
$\sum_{j=1}^k \lambda_j = 1$. A \emph{pure} decision rule puts all its weight on a
single action (e.g.\ ``always carry the umbrella''). A \emph{mixed} rule randomizes
(e.g.\ ``carry the umbrella with probability $0.7$'').

\medskip
\noindent
The \emph{reward of $\delta$ in environment $\omega$} is the average reward under
the randomization:
\[
  u(\delta, \omega) \;:=\; \sum_{j=1}^k \lambda_j\, u(a_j, \omega).
\]

\medskip
\noindent
\textbf{Risk vectors and the risk set.} To compare decision rules across all
environments simultaneously, we package the rewards into a single vector. The
\emph{risk vector} of a decision rule $\delta$ is
\[
  r(\delta) \;=\; \bigl(u(\delta, \omega_1),\; \ldots,\; u(\delta, \omega_n)\bigr)
  \;\in\; \mathbb{R}^n.
\]
Each coordinate records how well $\delta$ performs in one environment. In the umbrella
example, $r(a_1) = (8, 5)$ and $r(a_2) = (2, 10)$.

\medskip
\noindent
The \emph{risk set} $\mathcal{R}$ is the set of all risk vectors achievable by some
decision rule:
\[
  \mathcal{R} \;=\; \Bigl\{\sum_{j=1}^k \lambda_j\, r(a_j) \;:\;
  \lambda_j \geq 0,\; \sum_{j=1}^k \lambda_j = 1\Bigr\}.
\]
Geometrically, $\mathcal{R}$ is the \emph{convex hull} of the pure-action risk
vectors $\{r(a_1), \ldots, r(a_k)\}$: the set of all weighted averages of these
points.

\medskip
\noindent
\textbf{Admissibility (not being dominated).} A decision rule $\delta$ is
\emph{admissible} if there is no other rule $\delta'$ that does at least as well as
$\delta$ in \emph{every} environment and strictly better in at least one. If such a
$\delta'$ exists, we say $\delta'$ \emph{dominates} $\delta$, and any rational agent
should prefer $\delta'$ --- after all, switching from $\delta$ to $\delta'$ never
hurts and sometimes helps, regardless of which environment is the true one.

\medskip
\noindent
The \emph{Pareto frontier} $\mathcal{F} \subseteq \mathcal{R}$ is the set of risk
vectors of all admissible decision rules. A pure action $a_j$ is
\emph{Pareto-optimal} if $r(a_j) \in \mathcal{F}$.

\medskip
\noindent
\textbf{Bayesian expected utility.} A different approach to decision-making is to
assign a \emph{prior} $\pi = (\pi_1, \ldots, \pi_n)$ representing your beliefs about
how likely each environment is, where $\pi_i \geq 0$ and $\sum_i \pi_i = 1$. (For
example, $\pi = (0.3, 0.7)$ means you believe there is a $30\%$ chance of rain.) The
\emph{expected utility} of $\delta$ under $\pi$ is the weighted average reward:
\[
  \mathrm{EU}(\delta, \pi) \;:=\;
  \sum_{i=1}^n \pi_i\, u(\delta, \omega_i) \;=\; \pi \cdot r(\delta).
\]
A decision rule $\delta$ is \emph{Bayes-optimal} under $\pi$ if it achieves the
highest expected utility among all decision rules:
$\mathrm{EU}(\delta, \pi) \geq \mathrm{EU}(\delta', \pi)$ for all $\delta'$.

\medskip
\noindent
\textbf{The theorem.} These two approaches to rational decision-making ---
admissibility (``never use a dominated strategy'') and Bayesian expected utility
maximization (``assign beliefs and maximize average reward'') --- turn out to
characterize exactly the same set of decision rules:

\medskip
\noindent
\textbf{Theorem (Complete Class).} \emph{Every admissible decision rule is
Bayes-optimal under some prior $\pi$ with $\pi_i > 0$ for all $i$. Conversely, every
Bayes-optimal rule under such a prior is admissible.}

\medskip
\noindent
In other words: the decision rules that survive the ``no domination'' criterion are
precisely those that arise from maximizing expected utility under some set of beliefs
that doesn't rule out any environment entirely. This is significant because
admissibility is an extremely weak rationality requirement --- it says only that you
shouldn't use a strategy when a strictly better one is available --- yet it already
forces expected utility maximization.

\medskip
\noindent
\textbf{Part 3(a).} Show the following two facts:
\begin{itemize}
  \item[(i)] Any admissible decision rule $\delta = \sum_j \lambda_j\, a_j$ places
    zero weight on dominated pure actions: $\lambda_j = 0$ whenever $a_j$ is not
    Pareto-optimal. (In other words, the Pareto frontier $\mathcal{F}$ is contained
    in the convex hull of the Pareto-optimal pure actions alone.)\\
    \emph{Hint: If $\delta$ places positive weight on a dominated pure action $a_j$,
    replace $a_j$ with the action that dominates it. Does the resulting rule dominate
    $\delta$?}
  \item[(ii)] Every Bayes-optimal rule under a prior $\pi$ with $\pi_i > 0$ for all
    $i$ is admissible.\\
    \emph{Hint: If some $\delta'$ dominated $\delta$, compare their expected
    utilities. What does $\pi_i > 0$ ensure?}
\end{itemize}

\medskip
\noindent
\textbf{Part 3(b).} A \emph{face} $F$ of the Pareto frontier is a maximal convex
subset of $\mathcal{F}$ of the form
\[
  F \;=\; \Bigl\{\sum_{l=1}^p \mu_l\, r(a_{i_l}) \;:\;
  \mu_l \geq 0,\; \sum_{l=1}^p \mu_l = 1\Bigr\}
\]
for some subset of Pareto-optimal pure actions $\{a_{i_1}, \ldots, a_{i_p}\}$. Define
the \emph{difference vectors} $v_l := r(a_{i_l}) - r(a_{i_1})$ for
$l = 2, \ldots, p$, and let $H = \mathrm{span}\{v_2, \ldots, v_p\}$. Show that $H$
consists precisely of the directions along which one can move within $F$: that is, if
$r(\delta) \in F$, then $r(\delta) + h \in F$ for some $h$ only if $h \in H$.

\medskip
\noindent
\textbf{Part 3(c).} Let $\pi$ be a vector perpendicular to the subspace $H$ from
Part 3(b), normalized so that $\pi_i > 0$ for all $i$ and $\sum_i \pi_i = 1$. You may
assume that the entire risk set $\mathcal{R}$ lies on one side of the hyperplane
defined by $H$ (i.e.\ no point in $\mathcal{R}$ scores strictly higher under $\pi$
than the points on $F$). Show that:
\begin{itemize}
  \item[(i)] $\mathrm{EU}(\delta, \pi)$ takes the same value for all $\delta$ with
    $r(\delta) \in F$.
  \item[(ii)] Every rule on the face $F$ is Bayes-optimal under $\pi$.
\end{itemize}
Conclude the Complete Class Theorem: every admissible rule lies on some face of
$\mathcal{F}$, and the prior $\pi$ constructed from that face makes it Bayes-optimal.
\end{exercise}

\begin{solution}[ex:complete-class]
Higher reward is better; $\delta'$ \emph{dominates} $\delta$ if $r(\delta')\ge\rd$ coordinatewise with strict inequality in some coordinate. $\rd$ is linear in the mixing weights, $R=\operatorname{conv}\{r(a_1),\dots,r(a_k)\}$, and $\EU(\delta,\pi)=\pi\cdot\rd$.

\paragraph{Part 3(a)(i).} An admissible $\delta=\sum_j\lambda_j a_j$ puts $\lambda_j=0$ on every non-Pareto-optimal (dominated) $a_j$.

\soln Suppose $\lambda_j>0$ for a dominated $a_j$, and let $\delta'$ dominate it: $r(\delta')\ge r(a_j)$ with strict inequality in some coordinate $i_0$. Replace the weight on $a_j$ by $\delta'$, i.e.\ play $\delta'$ with the probability $\lambda_j$ formerly on $a_j$. This is a valid rule $\hat\delta$ with
\[
  r(\hat\delta)=\rd+\lambda_j\big(r(\delta')-r(a_j)\big)\ge\rd,
\]
and strict in coordinate $i_0$ (since $\lambda_j>0$). So $\hat\delta$ dominates $\delta$, contradicting admissibility. Hence $\lambda_j=0$, i.e.\ $F\subseteq\operatorname{conv}\{\text{Pareto-optimal pure actions}\}$. \qed

\paragraph{Part 3(a)(ii).} Every Bayes-optimal rule under a prior $\pi$ with $\pi_i>0$ for all $i$ is admissible.

\soln If $\delta'$ dominated such a $\delta$, then with $r(\delta')_{i_0}>\rd_{i_0}$,
\[
  \EU(\delta',\pi)-\EU(\delta,\pi)=\sum_i\pi_i\big(r(\delta')_i-\rd_i\big)\ge\pi_{i_0}\big(r(\delta')_{i_0}-\rd_{i_0}\big)>0,
\]
since every term is $\ge0$ and the $i_0$ term is $>0$. This contradicts Bayes-optimality, so $\delta$ is admissible. \qed

\paragraph{Part 3(b).} For a face $\mathcal F=\{\sum_{l=1}^p\mu_l\,r(a_{i_l}):\mu_l\ge0,\ \sum_l\mu_l=1\}$ with $v_l:=r(a_{i_l})-r(a_{i_1})$ and $H=\operatorname{span}\{v_2,\dots,v_p\}$: $H$ is exactly the set of directions tangent to $\mathcal F$.

\soln Any point of $\mathcal F$ equals $r(a_{i_1})+\sum_{l\ge2}\mu_l v_l$, so $\mathcal F\subseteq r(a_{i_1})+H$. Take $x,x+h\in\mathcal F$, say $x=\sum_l\mu_l r(a_{i_l})$ and $x+h=\sum_l\mu_l' r(a_{i_l})$ with $\sum_l\mu_l=\sum_l\mu_l'=1$. Then
\[
  h=\sum_l(\mu_l'-\mu_l)\,r(a_{i_l})=\sum_{l\ge2}(\mu_l'-\mu_l)\,v_l\in H,
\]
using $\sum_l(\mu_l'-\mu_l)=0$ to eliminate the $r(a_{i_1})$ term. Hence moving within $\mathcal F$ requires $h\in H$. Conversely, from any relative-interior point ($\mu_l>0$) every $h\in H$ is realized: $x+\varepsilon h\in\mathcal F$ for small $\varepsilon>0$. So the tangent directions of $\mathcal F$ are precisely $H$. \qed

\paragraph{Part 3(c).} Let $\pi\perp H$ be normalized so $\pi_i>0$ and $\sum_i\pi_i=1$, and assume no point of $R$ scores strictly higher under $\pi$ than the points of $\mathcal F$.

\soln \emph{(i)} For $x,x'\in\mathcal F$ we have $x-x'\in H$ by 3(b), so $\pi\cdot(x-x')=0$. Thus $\pi\cdot\rd$ equals a common value $c$ for all $\delta$ with $\rd\in\mathcal F$.

\emph{(ii)} Picture $\pi$ as the \emph{outward} normal of the supporting hyperplane $\{x:\pi\cdot x=c\}$: by assumption the whole risk set lies on the inner side, $\pi\cdot x\le c$ for all $x\in R$, touching the hyperplane exactly along the face $\mathcal F$. Fix any $\rd\in\mathcal F$ (so $\EU(\delta,\pi)=c$) and any other rule $\delta'$. The step from the face point $\rd$ to $r(\delta')\in R$ heads back into the risk set, i.e.\ against the outward normal, so its dot product with $\pi$ is non-positive:
\[
  \pi\cdot\big(r(\delta')-\rd\big)\le 0,\qquad\text{equivalently}\qquad \EU(\delta',\pi)\le\EU(\delta,\pi).
\]
Since this holds for every $\delta'$, each rule with $\rd\in\mathcal F$ maximizes expected utility, i.e.\ is Bayes-optimal under $\pi$.

\emph{Conclusion.} If $\delta$ is admissible then $\rd\in F$ lies on some face $\mathcal F$; the prior $\pi$ built from $\mathcal F$ has all $\pi_i>0$ and, by (ii), makes $\delta$ Bayes-optimal. Conversely, by 3(a)(ii) every Bayes-optimal rule under a strictly positive prior is admissible. The two classes coincide. \qed
\end{solution}

\begin{exercise}[The Do-Divergence Theorem]
\label{ex:do-divergence}
\noindent
\textbf{Motivation: optimization as steering.} A useful way to think about what it
means for an agent to be \emph{optimizing} is that it reliably steers the world into
a narrow set of outcomes --- outcomes that would be extremely unlikely to arise out of any random process . A thermostat keeps a room at 
20
∘
C
 despite varying weather; a
chess player steers toward checkmate despite the opponent's moves. In each case, the
actual outcome is concentrated in a small region of possibility space, whereas without
the agent's intervention, outcomes would be spread broadly.

\medskip
\noindent
This exercise makes that intuition precise in an information-theoretic setting. We
will show that an agent's ability to concentrate outcomes (``steer'') is bounded by
the amount of information the agent extracts from its observations. 

\medskip
\noindent
\textbf{Background: KL divergence.} Given two probability distributions $p$ and $q$
over the same set of outcomes, the \emph{Kullback--Leibler (KL) divergence} from $q$
to $p$ is
\[
  D_{\mathrm{KL}}(p \,\|\, q) \;:=\;
  \sum_x p(x) \log \frac{p(x)}{q(x)}.
\]
This quantity is always $\geq 0$, and equals $0$ only when $p = q$. It measures how
``different'' $p$ is from $q$, with a particular asymmetry: $D_{\mathrm{KL}}(p \| q)$
is large when $p$ places significant probability on outcomes where $q$ assigns very
little. In other words, it is large precisely when $p$ concentrates on outcomes that
would be \emph{surprising} under $q$. This can be seen as a signature of optimization:
the agent's policy makes certain outcomes likely that would be very unlikely under a
baseline policy.

\medskip
\noindent
\textbf{Background: mutual information.} The \emph{mutual information} between two
random variables $A$ and $O$ is
\[
  \mathrm{MI}(A;\, O) \;:=\;
  D_{\mathrm{KL}}\bigl(P[A, O] \,\big\|\, P[A]\,P[O]\bigr)
  \;=\; \sum_{a, o} P[a, o] \log \frac{P[a, o]}{P[a]\,P[o]}.
\]
This measures how much knowing $O$ tells you about $A$ (and vice versa). It is zero
when $A$ and $O$ are independent, and large when they are tightly coupled.

\medskip
\noindent
\textbf{Setup.} Consider an agent (the ``demon'') that observes some information $O$
about the world and then takes an action $A$ based on what it observed. The action and
observation together produce an outcome $X$. The joint distribution is
\[
  P[X, A, O] \;=\; P[X \mid A, O]\; P[A \mid O]\; P[O].
\]
The factor $P[A \mid O]$ encodes the demon's \emph{policy}: how it chooses actions as
a function of its observations.

\medskip
\noindent
Now consider a \emph{blind baseline}: the demon still acts, but ignores its
observations, choosing actions independently of $O$. We write the blind baseline
distribution as
\[
  P[X, A, O \mid do(A)] \;=\; P[X \mid A, O]\; P[A]\; P[O].
\]
The notation $do(A)$ means we have ``intervened'' on the action, replacing the
demon's observation-dependent policy $P[A \mid O]$ with the marginal $P[A]$ (the
overall frequency of each action, ignoring which observations prompted them). The
mechanism $P[X \mid A, O]$ by which actions and observations produce outcomes is
unchanged --- only the demon's strategy has been lobotomized.

\medskip
\noindent
\textbf{The theorem.} Prove the Do-Divergence Theorem:
\[
  D_{\mathrm{KL}}\bigl(P[X] \,\big\|\, P[X \mid do(A)]\bigr)
  \;\leq\; \mathrm{MI}(A;\, O).
\]
The left side measures how much the sighted demon's outcome distribution differs from
the blind baseline's. In the language of steering: it measures how much the demon has
concentrated outcomes into regions that would be unlikely without observation-dependent
action. The right side is the mutual information between actions and observations ---
how much the demon's actions depend on what it sees. The theorem says that \textbf{the
degree of steering is bounded by the information the demon uses}.

\medskip
\noindent
\emph{Useful fact:}
\begin{itemize}
  \item \textbf{Monotonicity of KL divergence.} For any two joint distributions
    $p(x,y)$ and $q(x,y)$, marginalizing out $y$ can only decrease KL divergence:
    $D_{\mathrm{KL}}(p(x) \,\|\, q(x)) \leq D_{\mathrm{KL}}(p(x,y) \,\|\, q(x,y))$.
    (Intuitively: forgetting information can only make two distributions look more
    similar, never less.)
\end{itemize}

\medskip
\noindent
\emph{Hint: Compute $D_{\mathrm{KL}}(P[X, A, O] \,\|\, P[X, A, O \mid do(A)])$
by expanding the log ratio using the factorizations above}

\medskip
\noindent
\textbf{Remark (Maxwell's demon and the thermodynamics of optimization).}

\medskip
\noindent
Maxwell's demon is a thought experiment in which a tiny intelligent being controls a
door between two halves of a box of gas. By observing each molecule's position and
selectively opening the door, the demon can sort all molecules to one side, creating a
highly ordered (low-entropy) state from an initially disordered one. In our notation:
the outcome $X$ is the final configuration of molecules, the observations $O$ are the
demon's measurements of molecular positions, and the actions $A$ are its door
openings.

\medskip
\noindent
Under the blind baseline (opening the door at random), molecules are roughly equally
likely to be on either side, so $P[X \mid do(A)]$ is spread broadly. If the demon
perfectly sorts all $n$ molecules to the right, $P[X]$ is concentrated on a single
configuration, and the KL divergence between these distributions is $n \log 2$ (i.e.\
$n$ bits). The theorem therefore says that perfectly sorting $n$ molecules requires
$\mathrm{MI}(A; O) \geq n$ bits: the demon must gather at least $n$ bits of
information about the molecules to reduce the gas's entropy by $n$ bits.
This illustrates the idea that any agent that steers a system into a narrow, unlikely region of outcome space (low entropy) must pay for this steering with mutual information. 
\end{exercise}

\begin{solution}[ex:do-divergence]
Write $\Pr[X,A,O]=\Pr[X\mid A,O]\Pr[A\mid O]\Pr[O]$ and the blind baseline $\Pr[X,A,O\mid\doo{A}]=\Pr[X\mid A,O]\Pr[A]\Pr[O]$.

\textbf{Claim.}
$\DKL\!\big(\Pr[X]\,\big\|\,\Pr[X\mid\doo{A}]\big)\le\MI(A;O)$.


\soln Compute the divergence between the full joints. The factors $\Pr[X\mid A,O]$ and $\Pr[O]$ cancel in the log-ratio:
\[
  \DKL\!\big(\Pr[X,A,O]\,\big\|\,\Pr[X,A,O\mid\doo{A}]\big)
  =\sum_{x,a,o}\Pr[x,a,o]\log\frac{\Pr[a\mid o]}{\Pr[a]}.
\]
The summand depends only on $a,o$; marginalizing $x$ and using $\Pr[a\mid o]/\Pr[a]=\Pr[a,o]/(\Pr[a]\Pr[o])$,
\[
  =\sum_{a,o}\Pr[a,o]\log\frac{\Pr[a,o]}{\Pr[a]\Pr[o]}=\MI(A;O).
\]
Marginalizing the two joints down to $X$ sends them to $\Pr[X]$ and $\Pr[X\mid\doo{A}]=\sum_{a,o}\Pr[X\mid a,o]\Pr[a]\Pr[o]$ respectively, so monotonicity of KL under marginalization gives
\[
  \DKL\!\big(\Pr[X]\,\big\|\,\Pr[X\mid\doo{A}]\big)\le\DKL\!\big(\Pr[X,A,O]\,\big\|\,\Pr[X,A,O\mid\doo{A}]\big)=\MI(A;O). \qedhere
\]
\end{solution}

\begin{exercise}[Channel Additivity]
\label{ex:channel-additivity}
\noindent
Consider two independent channels $X_1 \to Y_1$ and $X_2 \to Y_2$ with a fixed
joint channel
\[
  P[Y \mid X] \;=\; P[Y_1 \mid X_1]\, P[Y_2 \mid X_2].
\]
That is, output $Y_1$ depends only on input $X_1$, and $Y_2$ depends only on $X_2$;
the two channels do not interact. We are free to choose the input distribution $P[X]$
(which may correlate $X_1$ and $X_2$), and the joint distribution over everything is
then $P[Y, X] = P[Y \mid X]\, P[X]$. The goal is to maximize the mutual information
$\mathrm{MI}(X;\, Y)$, called the \emph{information throughput} of the channel.

\medskip
\noindent
\textbf{Part 5(a).} Show that for any input distribution $P[X]$,
\[
  \mathrm{MI}(X;\, Y) \;=\;
  \mathrm{MI}(X_1;\, Y_1) \;+\; \mathrm{MI}(X_2;\, Y_2) \;-\;
  \mathrm{MI}(Y_1;\, Y_2).
\]

\medskip
\noindent
\emph{Hint: Write each mutual information as a KL divergence between the joint and
the product of marginals, expand the log ratios using the channel factorization, and
collect terms.}

\medskip
\noindent
\textbf{Part 5(b).} Given any input distribution $P[X_1, X_2]$, define
\[
  Q[X_1, X_2] \;:=\; P[X_1]\, P[X_2],
\]
i.e.\ the product of the marginals. Show that
$\mathrm{MI}_Q(X;\, Y) \geq \mathrm{MI}_P(X;\, Y)$.

\medskip
\noindent
\emph{Hint: Under $Q$, the inputs are independent. What happens to
$\mathrm{MI}(Y_1;\, Y_2)$ when $X_1 \perp\!\!\!\perp X_2$, given the channel
factorization? Use Part 5(a).}

\medskip
\noindent
\textbf{Part 5(c).} Use Part 5(b) to conclude that there always exists a
throughput-maximizing input distribution under which $X_1 \perp\!\!\!\perp X_2$.

\medskip
\noindent
\textbf{Remark.} This result has a natural interpretation in terms of optimization and agency. Think of $X$ as the actions of an agent and $Y$ as the outcomes it cares about. The channel $P[Y \mid X]$ — how actions influence outcomes — is fixed by the environment and outside the agent's control; the agent only gets to choose its policy $P[X]$. The factorization condition on $P[Y \mid X]$ says that the environment is \emph{modular}: the two groups of outcomes $Y_1$ and $Y_2$ are each influenced only by their respective actions $X_1$ and $X_2$. Our theorem says that if the environment is modular in this sense, then the agent can always find an optimal policy that is modular in a corresponding sense --- specifically, the two groups of actions need not be coordinated at all and can be chosen independently.
\end{exercise}

\begin{solution}[ex:channel-additivity]
The channel factorizes as $P[Y\mid X]=P[Y_1\mid X_1]P[Y_2\mid X_2]$, and $P[Y,X]=P[Y\mid X]P[X]$.

\paragraph{Part 5(a).} $\MI(X;Y)=\MI(X_1;Y_1)+\MI(X_2;Y_2)-\MI(Y_1;Y_2)$.

\soln Using $P[x,y]=P[y\mid x]P[x]$ and the factorization,
\[
  \MI(X;Y)=\sum_{x,y}P[x,y]\log\frac{P[y\mid x]}{P[y]}
  =\sum_{x,y}P[x,y]\log\frac{P[y_1\mid x_1]P[y_2\mid x_2]}{P[y_1,y_2]}.
\]
The three terms on the right are, respectively,
\[
  \MI(X_1;Y_1)=\sum P[x,y]\log\tfrac{P[y_1\mid x_1]}{P[y_1]},\quad
  \MI(X_2;Y_2)=\sum P[x,y]\log\tfrac{P[y_2\mid x_2]}{P[y_2]},
\]
\[
  \MI(Y_1;Y_2)=\sum P[x,y]\log\tfrac{P[y_1,y_2]}{P[y_1]P[y_2]},
\]
each written over the full joint (the first depends only on $x_1,y_1$, etc.). Then
\[
  \MI(X_1;Y_1)+\MI(X_2;Y_2)-\MI(Y_1;Y_2)
  =\sum_{x,y}P[x,y]\log\frac{P[y_1\mid x_1]P[y_2\mid x_2]}{P[y_1,y_2]}=\MI(X;Y). \qed
\]

\paragraph{Part 5(b).} For $Q[X_1,X_2]:=P[X_1]P[X_2]$, $\MI_Q(X;Y)\ge\MI_P(X;Y)$.

\soln $Q$ keeps the marginals $Q[X_i]=P[X_i]$ and the channel, so $\MI_Q(X_i;Y_i)=\MI_P(X_i;Y_i)$ for $i=1,2$ ($\MI(X_i;Y_i)$ depends only on $P[X_i]$ and $P[Y_i\mid X_i]$). Under $Q$ the inputs are independent, so the outputs are too:
\[
  Q[y_1,y_2]=\sum_{x_1,x_2}Q[x_1]Q[x_2]P[y_1\mid x_1]P[y_2\mid x_2]=Q[y_1]\,Q[y_2],
\]
hence $\MI_Q(Y_1;Y_2)=0$. Applying 5(a) under each distribution,
\[
  \MI_Q(X;Y)=\MI_P(X_1;Y_1)+\MI_P(X_2;Y_2)
  =\MI_P(X;Y)+\MI_P(Y_1;Y_2)\ge\MI_P(X;Y),
\]
the last step using $\MI_P(Y_1;Y_2)\ge0$. \qed

\paragraph{Part 5(c).} A throughput-maximizing input with $X_1\perp X_2$ always exists.

\soln The input simplex is compact and $\MI(X;Y)$ continuous, so a maximizer $P^*$ exists. Let $Q^*:=P^*[X_1]P^*[X_2]$. By 5(b), $\MI_{Q^*}(X;Y)\ge\MI_{P^*}(X;Y)=\max$, so $Q^*$ is also a maximizer, and under $Q^*$ the inputs are independent. \qed
\end{solution}

\end{document}
